\documentclass{beamer}

\usepackage{cuhkbeamer}

\begin{document}

% ====================== Slide 1: TITLE SLIDE (REVISED) ======================
\begin{frame}
  \title{Infrastructure Economics and Management}
  \author{Chiu Yu KO}
  \institute{Low Altitude Economics}
  \date{Week 5}
  \titlepage
\end{frame}

% ====================== Slide 2: FORMAT & OBJECTIVES ======================
\begin{frame}[t]
\frametitle{Learning Objectives}
By the end of this class, you should be able to:
\begin{itemize}
\item \textbf{Compare} different vertiport archetypes and identify the main site-selection trade-offs.
\item \textbf{Estimate} how throughput, reliability, access time, and operating cost affect infrastructure viability.
\item \textbf{Calculate} breakeven activity and evaluate a staged ramp-up under uncertain demand.
\item \textbf{Assess} alternative methods for allocating scarce slots while protecting priority services.
\item \textbf{Evaluate} hub, distributed, and hybrid network designs for Hong Kong and the wider GBA.
\end{itemize}
\end{frame}

% ====================== Slide 3: ICEBREAKER ======================
\begin{frame}[t]
\frametitle{The Infrastructure Bottleneck}
\textbf{Imagine planning a new LAE service in Hong Kong.}
\begin{itemize}
\item Aircraft may be available, but the service also needs a safe operating site, power, communications, passenger or cargo access, maintenance support, and sufficient throughput.
\item A shortage of suitable infrastructure can constrain the entire service network.
\end{itemize}
\vspace{0.3cm}
\textbf{Managerial question:} Which resource is most likely to become the binding constraint: aircraft, landing capacity, charging, airspace access, or customer access?
\begin{itemize}
\item \textbf{Quick poll (2 min):} What is the most important infrastructure bottleneck for your Week 4 use case?
\end{itemize}
\end{frame}

% ====================== Session 1 Title Card ======================
\begin{frame}[c]
  \begin{center}
    \huge\color{cuhkpurple}\textbf{Session 1}\\
    \vspace{0.5cm}
    \Large Vertiport Types, Site Selection \& Throughput Economics
  \end{center}
\end{frame}

% ====================== Slide 4: CONCEPT PRIMER 1 (REVISED) ======================
\begin{frame}[t]
\frametitle{Concept Primer: Three Vertiport Archetypes}
For managerial analysis, sites can be grouped into three broad archetypes:
\begin{itemize}
\item \textbf{Building-based site:} Uses an existing rooftop or elevated structure. It may offer strong access but require structural, safety, power, and circulation upgrades.
\item \textbf{Integrated hub:} Provides multiple pads or stands together with charging, processing, maintenance, and passenger or cargo facilities.
\item \textbf{Ground-based site:} Uses available land and may be easier to expand, but can create longer access time to major customer locations.
\end{itemize}
\textit{These are classroom archetypes rather than an official Hong Kong classification. Actual design depends on aircraft, operations, safety requirements, and local approval.}
\end{frame}

% ====================== Slide 5: SITE SELECTION ECONOMICS ======================
\begin{frame}[t]
\frametitle{Site Selection: A Multi-Criteria Decision}
A suitable site must balance several considerations:
\begin{itemize}
\item \textbf{Safety and obstacle environment}: approach and departure paths, separation, downwash or outwash, emergency response, and surrounding uses.
\item \textbf{Customer access}: total door-to-door time, transfers, elevators, walking, and accessibility.
\item \textbf{Utilities and operations}: charging, communications, maintenance, storage, staffing, and recovery capability.
\item \textbf{Capacity and expandability}: pads, parking or stands, turnaround time, operating hours, and future growth.
\item \textbf{Cost and approval}: construction, adaptation, lease, environmental effects, community acceptance, and regulatory requirements.
\end{itemize}
\end{frame}


% ====================== NEW SLIDE: APPLIED TOOL - GENERALIZED SITE COST ======================
\begin{frame}[t]
\frametitle{Applied Tool: Comparing Sites}
\begin{block}{A Simple Annual Site-Cost Framework}
\[
\begin{split}
\text{Generalized Site Cost}={}&\text{Annualized Infrastructure Cost}+\text{Annual OPEX}\\
&+\text{Monetized Access Burden}+\text{Expected Disruption Cost}
\end{split}
\]
\end{block}
\begin{itemize}
\item Infrastructure cost includes construction and major site adaptation, net of residual value where relevant.
\item OPEX includes staffing, utilities, maintenance, insurance, communications, and recurring site costs.
\item Access burden converts additional passenger or cargo access time into a comparable value.
\item Disruption cost may reflect weather, maintenance, congestion, or other reliability risks.
\end{itemize}
\textit{Lower generalized cost is not automatically better if the site cannot meet safety, capacity, or service requirements.}
\end{frame}

% ====================== NEW SLIDE: BACK-OF-THE-ENVELOPE - SITE SELECTION ======================
\begin{frame}[t,shrink=12]
\frametitle{Back-of-the-Envelope: Comparing Two Site Archetypes}
\textbf{Assumptions:} 100 passenger-trips per day and value of access time = HK\$15 per minute.
\begin{columns}[T]
\begin{column}{0.48\textwidth}\begin{block}{Building-Based Site}
\begin{itemize}
\item Daily infrastructure allocation: HK\$5,000
\item Daily OPEX: HK\$8,000
\item Extra access time: 8 minutes
\item Access burden: $8\times15\times100=\text{HK\$12,000}$
\item \textbf{Generalized site cost: HK\$25,000/day}
\end{itemize}\end{block}\end{column}
\begin{column}{0.48\textwidth}\begin{block}{Ground-Based Site}
\begin{itemize}
\item Daily infrastructure allocation: HK\$1,000
\item Daily OPEX: HK\$2,000
\item Extra access time: 3 minutes
\item Access burden: $3\times15\times100=\text{HK\$4,500}$
\item \textbf{Generalized site cost: HK\$7,500/day}
\end{itemize}\end{block}\end{column}
\end{columns}
\textit{Takeaway: Under these assumptions, the ground site has lower generalized cost. Actual choice also depends on demand, safety, capacity, network value, and feasible access.}
\vfill{\tiny Illustrative assumptions for classroom analysis; not current site quotations.}
\end{frame}

% ====================== Slide 6: CONCEPT PRIMER 2 ======================
\begin{frame}[t]
\frametitle{Concept Primer: Throughput and Reliability}
\textbf{Throughput} is the number of aircraft movements, passengers, or missions that a facility can process during a defined period.
\begin{itemize}
\item It depends on pads, stands, turnaround, charging, passenger processing, staffing, airspace, and operating procedures.
\end{itemize}
\textbf{Reliability} is the facility's ability to provide the planned capacity when required.
\begin{itemize}
\item Downtime reduces usable capacity and may create delays, diversions, refunds, or lost demand.
\end{itemize}
\textbf{Managerial lesson:} Theoretical maximum throughput is less useful than reliable, serviceable capacity under realistic conditions.
\end{frame}

% ====================== Slide 7: STANDALONE FORMULA ======================
\begin{frame}[t]
\frametitle{Applied Tool: Breakeven Capacity Use}
\begin{block}{A Simple Facility Formula}
\[
\text{Breakeven Capacity Use}
=
\frac{\text{Fixed Daily Facility Cost}}
{\text{Reliable Daily Capacity}\times\text{Contribution per Movement}}
\]
\end{block}
\begin{itemize}
\item Contribution per movement is the facility fee minus variable cost associated with that movement.
\item Reliable capacity should reflect operating hours, turnaround, maintenance, weather, priority access, and other constraints.
\item The result is the share of reliable capacity that must be sold to cover fixed facility cost.
\item It is not equivalent to the percentage of hours in the day during which the facility must operate.
\end{itemize}
\end{frame}

% ====================== Slide 8: BACK-OF-ENVELOPE (BREAKEVEN) ======================
\begin{frame}[t,shrink=10]
\frametitle{Back-of-the-Envelope: Breakeven Capacity Use}
\textbf{Assumptions:} Reliable capacity = 100 movements/day; facility fee = HK\$200/movement; variable facility cost omitted.
\begin{columns}[T]
\begin{column}{0.48\textwidth}\begin{block}{Higher-Cost Site}
\begin{itemize}
\item Fixed cost: HK\$15,000/day
\item Maximum fee revenue: HK\$20,000/day
\item Breakeven use: $15{,}000/20{,}000$
\item \textbf{75\% of reliable capacity}
\end{itemize}\end{block}\end{column}
\begin{column}{0.48\textwidth}\begin{block}{Lower-Cost Site}
\begin{itemize}
\item Fixed cost: HK\$4,000/day
\item Maximum fee revenue: HK\$20,000/day
\item Breakeven use: $4{,}000/20{,}000$
\item \textbf{20\% of reliable capacity}
\end{itemize}\end{block}\end{column}
\end{columns}
\textit{Takeaway: A higher-cost site requires more activity, a higher contribution per movement, or complementary revenue.}
\vfill{\tiny Illustrative assumptions; excludes variable cost, time-of-day demand, and service mix.}
\end{frame}

% ====================== Slide 9: CONCEPT PRIMER 3 ======================
\begin{frame}[t]
\frametitle{Concept Primer: Staged Ramp-Up}
New infrastructure rarely reaches mature demand immediately.
\begin{itemize}
\item Early utilization may be limited by customer adoption, aircraft availability, approvals, operating experience, and network coverage.
\item A staged plan links investment to evidence and milestones rather than building the full facility at once.
\item Possible stages include a temporary or modular site, limited operating hours, selected users, and later expansion.
\item The financial plan must include working capital and liquidity for early operating losses and contingencies.
\end{itemize}
\textbf{Managerial question:} Which investment is reversible or expandable if demand develops differently from the forecast?
\end{frame}

% ====================== Slide 10: BACK-OF-ENVELOPE (RAMP-UP) ======================
\begin{frame}[t,shrink=10]
\frametitle{Back-of-the-Envelope: Funding the Ramp-Up}
\textbf{Assumptions:} Initial infrastructure investment = HK\$10 million; annual fixed operating cost = HK\$2 million. Other operating costs are omitted.
\begin{columns}[T]
\begin{column}{0.48\textwidth}\begin{block}{Mature-Demand Case}
\begin{itemize}
\item Facility use: 80\%
\item Annual revenue: HK\$4 million
\item Fixed operating cost: HK\$2 million
\item \textbf{Operating surplus before other costs: HK\$2 million}
\end{itemize}\end{block}\end{column}
\begin{column}{0.48\textwidth}\begin{block}{Year 1 Ramp-Up Case}
\begin{itemize}
\item Facility use: 20\%
\item Annual revenue: HK\$1 million
\item Fixed operating cost: HK\$2 million
\item \textbf{Operating shortfall before other costs: HK\$1 million}
\end{itemize}\end{block}\end{column}
\end{columns}
\textit{Takeaway: The financing plan must cover the expected ramp-up shortfall plus a contingency buffer.}
\vfill{\tiny Illustrative assumptions; excludes variable cost, financing, depreciation, tax, and working-capital timing.}
\end{frame}

% ====================== Slide 11: STANDALONE FORMULA ======================
\begin{frame}[t]
\frametitle{Applied Tool: A Reliability Scenario}
\begin{block}{A Simple Annual Return Screen}
\[
\text{Simple Return}
=
\frac{\text{Adjusted Revenue}-\text{Adjusted Operating Cost}}
{\text{Initial Investment}}
\]
\end{block}
\begin{itemize}
\item Adjusted revenue should reflect serviceable capacity, demand, cancellations, and pricing under the scenario.
\item Some costs fall during downtime, while fixed cost, recovery cost, refunds, and service credits may remain or increase.
\item Weather should be modelled from route- and aircraft-specific operating assumptions rather than a universal annual percentage.
\item This simple return is not a full ROI or NPV analysis because it ignores timing, financing, tax, and residual value.
\end{itemize}
\end{frame}

% ====================== Slide 12: BACK-OF-ENVELOPE (WEATHER HAIRCUT) ======================
\begin{frame}[t,shrink=12]
\frametitle{Back-of-the-Envelope: A Downtime Scenario}
\textbf{Assumptions:} Initial investment = HK\$10 million; annual revenue without disruption = HK\$4 million; annual cost = HK\$2.5 million.
\begin{columns}[T]
\begin{column}{0.48\textwidth}\begin{block}{No-Downtime Scenario}
\begin{itemize}
\item Operating surplus: HK\$1.5 million
\item Simple return: $1.5/10$
\item \textbf{15\%}
\end{itemize}\end{block}\end{column}
\begin{column}{0.48\textwidth}\begin{block}{20\% Revenue Reduction}
\begin{itemize}
\item Adjusted revenue: HK\$3.2 million
\item Assumed cost: HK\$2.5 million
\item Operating surplus: HK\$0.7 million
\item \textbf{Simple return: 7\%}
\end{itemize}\end{block}\end{column}
\end{columns}
\textit{Takeaway: The assumed disruption materially affects returns. A better model would separate capacity loss, demand loss, variable-cost savings, and recovery cost.}
\vfill{\tiny The 20\% reduction is an illustrative stress scenario, not a claim about Hong Kong annual weather downtime.}
\end{frame}

% ====================== Slide 13: SESSION 2 TITLE ======================
\begin{frame}[c]
  \begin{center}
    \huge\color{cuhkpurple}\textbf{Session 2}\\
    \vspace{0.5cm}
    \Large CAPEX, OPEX and the Maintenance Money-Pit
  \end{center}
\end{frame}

% ====================== Slide 14: CONCEPT PRIMER 4 ======================
\begin{frame}[t]
\frametitle{Concept Primer: Infrastructure Investment}
Vertiport investment may include:
\begin{itemize}
\item site acquisition or adaptation and structural works;
\item touchdown, lift-off, safety, parking, and passenger or cargo areas;
\item charging, energy storage, communications, lighting, security, and fire protection;
\item access, circulation, accessibility, and integration with ground transport; and
\item design, engineering, testing, approval, and contingency allowances.
\end{itemize}
\textit{Requirements depend on aircraft characteristics, operating concept, site conditions, and the applicable regulatory framework.}
\end{frame}

% ====================== Slide 15: CONCEPT PRIMER 5 ======================
\begin{frame}[t]
\frametitle{Concept Primer: Infrastructure OPEX}
Recurring facility cost may include:
\begin{itemize}
\item staffing, security, emergency response, and passenger or cargo processing;
\item electricity, communications, software, monitoring, and data services;
\item inspection, preventive maintenance, repairs, cleaning, and replacement reserves;
\item lease payments, insurance, professional services, and recurring compliance activities; and
\item disruption recovery, customer support, and service-level obligations.
\end{itemize}
\textbf{Managerial lesson:} Do not assume electricity is the dominant OPEX item. The cost mix depends on the site and operating model.
\end{frame}

% ====================== Slide 16: STANDALONE FORMULA ======================
\begin{frame}[t]
\frametitle{Applied Tool: Daily Infrastructure Cost}
\begin{block}{A Simple Economic-Cost Approximation}
\[
\text{Daily Infrastructure Cost}
=
\frac{\text{CAPEX}-\text{Expected Residual Value}}{\text{Useful Operating Days}}
+\text{Daily OPEX}
\]
\end{block}
\begin{itemize}
\item The formula allocates net investment cost over the expected operating life.
\item It is useful for unit-cost screening but is not identical to daily cash flow, accounting depreciation, mortgage payment, ROI, or NPV.
\item A complete investment model should also consider financing, replacement cycles, tax, residual value, and the timing of cash flows.
\end{itemize}
\end{frame}

% ====================== Slide 17: BACK-OF-ENVELOPE (DAILY BURDEN) ======================
\begin{frame}[t,shrink=10]
\frametitle{Back-of-the-Envelope: Daily Infrastructure Cost}
\textbf{Assumption:} Ten-year life and no residual value.
\begin{columns}[T]
\begin{column}{0.48\textwidth}\begin{block}{Higher-Investment Site}
\begin{itemize}
\item CAPEX: HK\$18.25 million
\item Daily allocation: HK\$5,000
\item Daily OPEX: HK\$10,000
\item \textbf{Daily economic cost: HK\$15,000}
\end{itemize}\end{block}\end{column}
\begin{column}{0.48\textwidth}\begin{block}{Lower-Investment Site}
\begin{itemize}
\item CAPEX: HK\$3.65 million
\item Daily allocation: HK\$1,000
\item Daily OPEX: HK\$3,000
\item \textbf{Daily economic cost: HK\$4,000}
\end{itemize}\end{block}\end{column}
\end{columns}
\textit{Takeaway: A higher-cost site must create sufficient additional demand, price, reliability, network value, or complementary revenue to justify the difference.}
\vfill{\tiny Illustrative assumptions; excludes financing, tax, residual value, and replacement cycles.}
\end{frame}

% ====================== Slide 18: SESSION 3 TITLE ======================
\begin{frame}[c]
  \begin{center}
    \huge\color{cuhkpurple}\textbf{Session 3}\\
    \vspace{0.5cm}
    \Large Slot Allocation, EMS Priority \& GBA Scale
  \end{center}
\end{frame}

% ====================== Slide 19: CONCEPT PRIMER 6 ======================
\begin{frame}[t]
\frametitle{Concept Primer: Allocating Scarce Slots}
When requests exceed safe and reliable capacity, managers need a transparent allocation rule.
\begin{itemize}
\item \textbf{Administrative allocation:} fixed schedules, permits, historic rights, rotation, or first-come-first-served.
\item \textbf{Market-based allocation:} posted peak pricing or auctions where appropriate.
\item \textbf{Priority rules:} emergency, public-service, safety, or other protected operations may receive priority under predefined procedures.
\item \textbf{Operational controls:} reservations, buffers, delay management, diversions, and cancellation rules.
\end{itemize}
\textbf{Managerial objective:} Balance safety, efficiency, predictability, access, competition, and public-service obligations.
\end{frame}

% ====================== Slide 20: STANDALONE FORMULA ======================
\begin{frame}[t]
\frametitle{Applied Tool: A Peak-Pricing Rule}
\begin{block}{An Illustrative Posted-Price Rule}
\[
\text{Slot Fee}
=
\text{Base Fee}
+
\text{Congestion Charge}\left(\frac{\text{Requests}}{\text{Available Slots}}\right)
\]
\end{block}
\begin{itemize}
\item The exact pricing function is a managerial design choice, not a standard regulatory formula.
\item Peak charges may shift flexible demand, but only if operators can change timing, route, or facility.
\item Pricing must be coordinated with safety limits, priority rules, competition policy, transparency, and access obligations.
\item Revenue maximization should not be the only objective of scarce infrastructure allocation.
\end{itemize}
\end{frame}

% ====================== Slide 21: BACK-OF-ENVELOPE (CONGESTION PRICING) ======================
\begin{frame}[t,shrink=10]
\frametitle{Back-of-the-Envelope: An Illustrative Peak Fee}
\textbf{Assumptions:} Base fee = HK\$100 and an illustrative congestion parameter = HK\$200.
\begin{columns}[T]
\begin{column}{0.48\textwidth}\begin{block}{Two Requests, Two Slots}
\begin{itemize}
\item Demand-capacity ratio: $2/2=1$
\item Fee: $100+1(200)$
\item \textbf{HK\$300}
\end{itemize}\end{block}\end{column}
\begin{column}{0.48\textwidth}\begin{block}{Eight Requests, Two Slots}
\begin{itemize}
\item Demand-capacity ratio: $8/2=4$
\item Fee: $100+4(200)$
\item \textbf{HK\$900}
\end{itemize}\end{block}\end{column}
\end{columns}
\textit{Takeaway: The higher fee may shift flexible trips, but customer response must be estimated. Price alone does not identify which missions create the greatest social value.}
\vfill{\tiny Illustrative pricing rule and assumptions; not a recommended tariff.}
\end{frame}

% ====================== Slide 22: HK VS GBA COMPARISON (REVISED) ======================
\begin{frame}[t]
\frametitle{Hong Kong and the GBA: Infrastructure Possibilities}
\begin{center}\small
\begin{tabular}{l|cc}
\textbf{Dimension} & \textbf{Hong Kong possibility} & \textbf{Wider GBA possibility}\\
\hline
Site environment & Dense and constrained & More varied\\
Customer access & Strong value in central locations & Varies by city and corridor\\
Infrastructure role & Integration and specialized services & Larger network and production scale\\
Likely design & Hub, distributed, or hybrid & Hub, distributed, or hybrid\\
Key question & Value created per scarce site & Network coverage and volume\\
\end{tabular}
\end{center}
\vspace{0.3cm}
\textit{Network design should be based on route-specific demand, site feasibility, aircraft, regulation, and access rather than a predetermined Hong Kong or GBA model.}
\end{frame}
% ====================== Slide 23: CONCEPT PRIMER 7 ======================
\begin{frame}[t]
\frametitle{Concept Primer: Network Design}
Should the network concentrate activity in a few hubs or distribute activity across many sites?
\begin{itemize}
\item \textbf{Concentrated hub network:} May support shared facilities, high throughput, connections, and specialized services, but can increase access time and vulnerability to disruption.
\item \textbf{Distributed network:} May improve customer access and resilience, but can duplicate fixed cost and reduce activity at each site.
\item \textbf{Hybrid network:} Combines major hubs with smaller feeder or specialized sites.
\end{itemize}
\textbf{Managerial lesson:} There is no universal ``Hong Kong strategy'' or ``GBA strategy.'' The best configuration depends on demand, geography, capacity, and service purpose.
\end{frame}

% ====================== NEW SLIDE: APPLIED TOOL - NETWORK DESIGN COST ======================
\begin{frame}[t]
\frametitle{Applied Tool: Comparing Network Designs}
\begin{block}{Generalized Cost per Served Trip}
\[
\text{Generalized Cost per Trip}
=
\frac{\text{Network Fixed Cost}+\text{Network OPEX}+\text{Disruption Cost}}{\text{Served Trips}}
+\text{Access Burden per Trip}
\]
\end{block}
\begin{itemize}
\item Concentration may lower facility cost per trip but increase passenger or cargo access burden.
\item Distribution may improve access and resilience while increasing duplicated cost.
\item Served trips, not unconstrained forecast demand, should be used in the denominator.
\item Compare service quality, reliability, expansion options, and network connections as well as cost.
\end{itemize}
\end{frame}

% ====================== NEW SLIDE: BACK-OF-THE-ENVELOPE (NETWORK DESIGN) ======================
\begin{frame}[t,shrink=12]
\frametitle{Back-of-the-Envelope: Concentrated vs Distributed Network}
\textbf{Assumptions:} 4,000 served passenger-trips/day and value of access time = HK\$15/minute.
\begin{columns}[T]
\begin{column}{0.48\textwidth}\begin{block}{Concentrated Network}
\begin{itemize}
\item Daily network cost: HK\$15,000
\item Cost per trip: HK\$3.75
\item Access burden: $8\times15=\text{HK\$120/trip}$
\item \textbf{Generalized cost: HK\$123.75/trip}
\end{itemize}\end{block}\end{column}
\begin{column}{0.48\textwidth}\begin{block}{Distributed Network}
\begin{itemize}
\item Daily network cost: HK\$45,000
\item Cost per trip: HK\$11.25
\item Access burden: $2\times15=\text{HK\$30/trip}$
\item \textbf{Generalized cost: HK\$41.25/trip}
\end{itemize}\end{block}\end{column}
\end{columns}
\textit{Takeaway: The distributed design wins under these assumed volumes and access times. The result may reverse if demand, reliability, feasible sites, or network cost changes.}
\vfill{\tiny Illustrative assumptions; 4,000 refers to passenger-trips, not flights.}
\end{frame}

% ====================== Slide 24: EXERCISE INSTRUCTIONS ======================
\begin{frame}[t]
\frametitle{In-Class Exercise: Infrastructure Viability Workshop}
\textbf{Activity: Stress-Testing a Two-Pad Facility}
\begin{itemize}
\item \textbf{Format:} Groups of 4--5
\item \textbf{Base task:} Estimate reliable daily capacity, contribution per movement, breakeven capacity use, and simple annual return.
\item \textbf{Scenario 1:} Reduce serviceable capacity to reflect weather or maintenance downtime.
\item \textbf{Scenario 2:} Reserve five daily slots for priority public-service use.
\item \textbf{Decision:} Recommend a fee, availability payment, subsidy, or cost reduction that keeps the project viable without violating capacity or priority rules.
\item \textbf{Timing:} 15 minutes plus 5-minute discussion.
\end{itemize}
\end{frame}

% ====================== Slide 25: KEY TAKEAWAYS ======================
\begin{frame}[t]
\frametitle{Key Takeaways}
\begin{itemize}
\item Infrastructure strategy balances safety, customer access, utilities, reliable capacity, cost, expandability, and community impact.
\item Breakeven should use reliable capacity and contribution per movement, not theoretical maximum activity or gross fees alone.
\item Staged investment and liquidity planning reduce the risk of building ahead of demand.
\item Weather and maintenance should be modelled through explicit capacity, demand, cost, and recovery assumptions.
\item Scarce slots require transparent allocation and priority rules; peak pricing is only one possible mechanism.
\item Hub, distributed, and hybrid networks each involve trade-offs between facility cost, access, resilience, and connectivity.
\end{itemize}
\end{frame}

% ====================== Slide 26: ASSIGNMENT ======================
\begin{frame}[t]
\frametitle{Week 5 Assignment Briefing}
\textbf{Infrastructure Viability Model and Memo (maximum 1,200 words plus spreadsheet)}
\begin{itemize}
\item \textbf{Due:} Beginning of Week 6
\item \textbf{Task:} Evaluate a building-based, ground-based, or integrated-hub archetype for your selected use case.
\item \textbf{Requirements:}
\begin{enumerate}
\item Explain the site-selection criteria and define reliable daily capacity.
\item Separate infrastructure investment, fixed OPEX, variable facility cost, and access burden.
\item Calculate breakeven capacity use and a simple annual return.
\item Model a staged ramp-up and one weather or maintenance scenario.
\item Propose a slot-allocation rule that protects priority services.
\item Compare the selected design with one alternative and recommend whether to build, stage, relocate, or reject the project.
\end{enumerate}
\end{itemize}
\end{frame}

% ====================== Slide 27: THANK YOU ======================
\begin{frame}[t]
  \frametitle{Thank You + Q\&A}
  
  \begin{center}
      \textbf{Next Week Preview:} \\
      \Large Public-Private Partnerships: Who pays for the chargers?
  \end{center}
\end{frame}

\end{document}